\documentclass[11pt]{article}
\usepackage{amsmath, amssymb, amscd, amsthm, amsfonts}
\usepackage{graphicx}
\usepackage{hyperref}
\usepackage{subcaption}
\usepackage{url}

\oddsidemargin 0pt
\evensidemargin 0pt
\marginparwidth 40pt
\marginparsep 10pt
\topmargin -20pt
\headsep 10pt
\textheight 8.7in
\textwidth 6.65in
\linespread{1.2}

\title{\textbf{From IDE to Edge: A Cloud-Native Multi-Agent Framework for Automated Edge AI Deployment on STM32}}
\author{\textbf{Supervisors}\\Prof. Gianvito Urgese  \\ Andrea Pignata \\ Giuseppe Fanuli \and \textbf{Candidate} \\ Michele \textsc{Russo}}
\date{March, 2026\\Politecnico di Torino}

\begin{document}

\maketitle

\section*{Summary}\label{section:summary}

\subsection*{Background}

Deploying AI models to constrained embedded systems remains a multi-step, error-prone process that demands expertise in both machine learning and embedded systems engineering. On STM32 microcontrollers, a practitioner must configure the peripheral hardware with STM32CubeMX, quantize and validate the model with STEdgeAI, fine-tune or compress the model to fit within strict Flash and RAM budgets, and then manually patch the generated C code to integrate the inference engine into the firmware project. Each step is handled by a separate tool with its own CLI, configuration format, and failure mode. Connecting them together is typically done by hand-crafted shell scripts or manual copy-paste operations. This friction causes long iteration cycles and raises the barrier to entry for developers who are not simultaneously deep experts in TensorFlow, embedded C, and ARM toolchains.

Large Language Models (LLMs) offer a natural interface for this kind of orchestration: they can parse free-text intent, select parameters, invoke tools, and interpret error messages. However, a single LLM call is not reliable enough for multi-step engineering tasks. Isolated queries to a chat interface do not maintain state across tool invocations, cannot route to different expertise levels for different subtasks, and offer no mechanism for a human engineer to review or override intermediate decisions before they cascade into irreversible states.

The emerging multi-agent paradigm addresses these limitations by decomposing the pipeline into a directed graph of specialized nodes, each responsible for one phase of the lifecycle. State is passed explicitly between nodes via a typed schema, enabling conditional routing, checkpointing, and supervised handoffs to a human operator at defined decision points. Applied to Edge AI development for STM32, this paradigm enables a system that can take a high-level natural language request (e.g., \textit{``Deploy a gesture classifier on STM32H7 with high compression''}) and produce a verified, compilable firmware project with minimal manual intervention.

\subsection*{Thesis Objectives and Contributions}

This thesis designs, implements, and evaluates a multi-agent MLOps orchestration framework for automating the full Edge AI deployment cycle on STM32 microcontrollers. The system is built on \textbf{LangGraph} and manages seven specialized workflows as subgraphs connected through a shared typed state and a \textbf{Redis}-backed checkpointing layer. Human-in-the-Loop (HITL) gates, implemented via LangGraph's \texttt{interrupt()} primitive, allow engineers to approve decisions, redirect searches, or override model choices at defined breakpoints without losing execution context.

The inference backend evolved from local \textbf{Ollama} instances to a centralized \textbf{NVIDIA Triton Inference Server} hosted on the \textbf{inNuCE Lab Heterogeneous Prototyping Platform (HPP)} cluster (NVIDIA A40, 45~GB VRAM), which eliminated the VRAM saturation crashes that occurred when multiple users loaded different models concurrently on a 16~GB development GPU.

\bigskip

The contributions of this thesis are as follows:

\begin{enumerate}

    \item A \textbf{multi-agent Edge AI orchestration framework} built on LangGraph that manages the complete STM32 deployment pipeline, from CubeMX firmware project generation and CMSIS-NN model analysis to fine-tuning, NNI hyperparameter optimization, and binary integration. The system uses a headless CubeMX automation strategy (\texttt{xvfb-run} virtual framebuffer) to bypass the GUI bottleneck in non-interactive environments.

    \item A \textbf{hardware-aware model customization pipeline} that applies heuristic anti-overfitting rules (layer freezing, dropout injection, epoch capping) on a fast path and a deterministic NNI template generator on the optimization path. A class-mismatch rescue mechanism automatically replaces the final classification layer when the pre-trained model output dimension conflicts with the target dataset.

    \item A \textbf{scalable inference architecture} progressing through three deployment scenarios (local workstation, per-user Docker containers with Ollama, centralized Triton on HPP) with quantitative stress tests at 5, 10, and 15 concurrent users. At 15 users, Triton remained stable at 80\% VRAM occupancy with GPU compute near 100\%, demonstrating that the 1~Gbps network connection was not a throughput bottleneck.

    \item A \textbf{developer productivity evaluation} based on KLM task decomposition, measuring a median 9$\times$ end-to-end speedup across six deployment tasks (with peaks of 40$\times$ for CubeMX project setup and 30$\times$ for NNI experiment generation) compared to equivalent manual procedures on official documentation baselines.

    \item A \textbf{RAG quality assessment} using DeepEval's LLM-as-Judge framework across five representative queries, measuring Faithfulness, Contextual Relevancy, Hallucination, and Answer Relevancy. Results show an average Contextual Relevancy of 0.75 and a Hallucination rate of 0.14, with a clearly measurable Faithfulness-Utility trade-off between procedural and exploratory query types.

\end{enumerate}

\end{document}